\noindent\textbf{Can a model read a printed sentence as an image and write a translation as a stream of glyphs on one consumer GPU? We trained two matched English--Chinese writers on 39,008 bilingual pairs and 352 additional comparisons of a reference with the initial writer's own image. On 96 repeated correction prompts, whole-answer preference training raised complete-meaning credits from zero to 65. Likelihood training earned 61. On other rehearsed passages the scores were eight and seven. On four unseen books the preference writer earned one credit in 96 attempts, for a figure label that also earned credit before training. A second experiment gave the original writer access to 497,775 pictured translation pairs during generation. Related passages earned two credits on 96 requests from four further books. Random passages and the original writer each earned one. Both approaches improved some glyph decisions. Neither produced reliable translations of new prose.}\par
\begin{multicols}{2}
\section*{Correcting an emitted glyph sequence}
The editor reads the pictured source and a rough visual translation, then generates glyphs or copies them from the draft.\cite{pointer} Both arms rehearsed the same targets. One also learned to score the entire reference above a previously emitted error, using mean log probability per glyph. Two bilingual graders screened the errors. An individual recheck reduced 400 proposed comparisons to 352, with some omissions still accepted.
\end{multicols}
\displayfigure{Visual_sequence_preference.pdf}{Learning a sentence and writing on a new page}{Top, the first repeated training example credited only to the preference writer. Middle, both the preference writer and its parent preserved figure number 13. Bottom, an unseen request for a summary ends in repetitions of \emph{for us}. All eight strips are original input or output pixels.}
\begin{multicols}{2}
\section*{Testing transfer}
The text comparator earned 92/96 book credits, supplied references 71 and unrelated foils zero. The preference writer preserved both polarities in 197/236 controlled pairs, against 225 before further training. A separate CPU run reproduced all 24 sampled output images and stops.
\section*{Where the writer loses meaning}
On the books, every neural writer assigned its free output a higher mean glyph probability than the reference in all 96 cases. With the correct prefix provided, related bilingual passages supplied the right next cell on 187 of 1,274 parent errors, versus 54 for random passages. We next tested the same evidence after the writer had emitted its own preceding cells. \repo{references/visual_sequence_preference_result.md}{Methods, scores and raw images}.
\end{multicols}
